\chapter{Conclusion and Future Work}
\label{ch:conclusion}

\section{Answer to the Main Research Question}
Combining camera and sensor predictions modestly improved completed-action failure
ranking over logistic sensing in REASSEMBLE. Logistic stacking achieved failure AP
0.4473, compared with 0.4294 for sensors and 0.4457 for equal probability averaging.
Equal averaging scored similarly, and its difference from stacking remained unresolved.
The result supports a conditional benefit from combining the available observations,
without establishing that a learned combiner was necessary.

The stacking gain over sensors did not recur in MSM or ViFailback. HANDOVER tested
averaging on an official split and also favored its sensor component, while unavailable
participant identities prevented the planned learned-fusion comparison. Across held-out
gear experiments, fusion of three sensor groups showed no clear improvement over the
training mean when predicting the photographic proxy, whose relationship to physical
tooth damage remained unvalidated.
Chapter~\ref{ch:discussion} gives the detailed contrasts and their study-specific
boundaries.

Together, the studies support evaluating each stream and simple fusion baselines before
adding adaptive weighting. Consistent adaptive benefit under degraded but present inputs
remained unresolved. The evidence concerns condition estimation and retrospective
assessment of completed events. It does not establish prospective weld-quality
decisions, which require an independently verified outcome and observations available
at a prespecified decision time.

\section{Supported Contributions}
\textbf{SQ1 --- Representation comparison.} The thesis compared statistical sensor
summaries, temporal encoders and visual representations under explicit task and fitting
conditions. It identified useful combinations of sensor representations and promising
visual pooling choices, while separating those findings from claims about learned
temporal order or universal encoder superiority.

\textbf{SQ2 --- Incremental fusion value.} The thesis estimated what fusion added over
single-input predictors and equal averaging on common assessment observations. The
contrasting assembly and external results show why a fusion claim depends on its
comparator, target and evaluation protocol.

\textbf{SQ3 --- Reliability under input failure.} The experiments distinguished
programmed fallback for a declared missing stream from adaptation to degraded inputs
that remained present. Matched-exposure comparisons tested whether learned weighting
added value beyond static fusion, leaving consistent adaptive benefit unresolved.

\textbf{SQ4 --- Interpretation beyond accuracy.} The evaluation combined ranking,
probability error, class-specific decisions, operating thresholds and grouped
uncertainty. This exposed cases in which better ranking did not imply better
probabilities or useful fixed-threshold decisions, and where changing the grouping
changed the uncertainty attached to the same predictions.

\section{Principal Limitations}
Each dataset imposes a different evidential limit. PHM lacks an independent
physical-damage reference, while native outcome definitions and ambiguous robotic events
constrain what the classification results mean. Incomplete identifiers for tools,
objects, sessions or participants limit claims of independence beyond the groups
actually checked.

Several follow-ups reused assessment populations examined during exploratory
development. Their intervals condition on saved fits and do not capture uncertainty from
the full selection process. Differences in calibration, training budgets and corruption
exposure also limit attribution of some contrasts to architecture. Preprocessing before
the separate four-arm PHM evaluation further compromises that diagnostic's isolation of
held-out observations.

Incomplete PHM historical training-source lineage and unavailable ViFailback per-fit
histories limit independent reconstruction beyond the retained configurations and
results. ViFailback target-CNN adaptation was assessed on one outer fold, with
inconclusive matched effects, so its consistency across task folds and seeds remains
unknown. Robustness to corrupted ViFailback inputs was not measured.

\section{Priorities for Future Research}
An independent welding study is the most direct next test of the original question. It
should establish weld quality through independent inspection, specify the prediction
horizon before examining outcomes, and reserve complete new welding sessions for
assessment. Camera and process measurements should cover only the information available
at that decision time. Warning time, false alarms and end-to-end delay would then
connect predictive performance to prospective use. For PHM, independent damage
measurements linked to known teeth would separate physical validity from agreement with
the photographic proxy.

For REASSEMBLE, a prespecified comparison on untouched recordings should test
last-quarter visual pooling, the same-sensor hybrid and anchored fusion separately and
in combination. Matched controls should preserve the relevant fitting, calibration and
corruption-exposure conditions to estimate their individual and joint effects. These
choices were promising in separate exploratory studies; their joint benefit has not
been established.

For ViFailback, extending matched fine-tuning across the remaining task folds and seeds
would test whether target-CNN adaptation changes visual ranking and fusion performance
consistently beyond the currently assessed fold. Missing, degraded and temporally
shifted inputs require a separate corruption study with its own scope and resource
estimate.

Secondary tests could examine HGB-based stacking in MSM and compare uncalibrated with
rank-preserving calibrated branch outputs in ViFailback, refitting F2 on each
corresponding training-side prediction bank. These are proposed comparisons with open
outcomes. Across them, fixed protocols, strong single-input controls and prespecified
tolerances for clean-performance loss would make any benefit easier to interpret.
